% problem_id: S001-lemma-cohomology-perfect-complex
% source_id: S001 (Stacks Project)
% source_file: quot.tex
% statement_locator: quot.tex lines 432--453
% proof_locator: quot.tex lines 455--581
% env: lemma
% id_note: label 在本来源内唯一
% pairing: tight (verified)
% status: complete (statement + author proof verbatim + reason + locators)
%
\begin{lemma}
\label{lemma-cohomology-perfect-complex}
Let $S$ be a scheme. Let $B$ be an algebraic space over $S$.
Let $K$ be a pseudo-coherent object of $D(\mathcal{O}_B)$.
\begin{enumerate}
\item If for all $g : T \to B$ in $(\Sch/B)$ the cohomology sheaf
$H^{-1}(Lg^*K)$ is zero, then the functor
$$
(\Sch/B)^{opp} \longrightarrow \textit{Sets},\quad
(g : T \to B) \longmapsto H^0(T, H^0(Lg^*K))
$$
is an algebraic space affine and of finite presentation over $B$.
\item If for all $g : T \to B$ in $(\Sch/B)$ the cohomology sheaves
$H^i(Lg^*K)$ are zero for $i < 0$, then $K$ is perfect,
$K$ locally has tor amplitude in $[0, b]$, and the functor
$$
(\Sch/B)^{opp} \longrightarrow \textit{Sets},\quad
(g : T \to B) \longmapsto H^0(T, Lg^*K)
$$
is an algebraic space affine and of finite presentation over $B$.
\end{enumerate}
\end{lemma}

\begin{proof}
Under the assumptions of (2) we have $H^0(T, Lg^*K) = H^0(T, H^0(Lg^*K))$.
Let us prove that the rule $T \mapsto H^0(T, H^0(Lg^*K))$ satisfies the
sheaf property for the fppf topology. To do this assume we have an
fppf covering $\{h_i : T_i \to T\}$ of a scheme $g : T \to B$ over $B$.
Set $g_i = g \circ h_i$. Note that since $h_i$ is flat, we have
$Lh_i^* = h_i^*$ and $h_i^*$ commutes with taking cohomology. Hence
$$
H^0(T_i, H^0(Lg_i^*K)) =
H^0(T_i, H^0(h_i^*Lg^*K)) =
H^0(T, h_i^*H^0(Lg^*K))
$$
Similarly for the pullback to $T_i \times_T T_j$.
Since $Lg^*K$ is a pseudo-coherent complex on $T$
(Cohomology on Sites, Lemma
\ref{sites-cohomology-lemma-pseudo-coherent-pullback})
the cohomology sheaf $\mathcal{F} = H^0(Lg^*K)$ is quasi-coherent
(Derived Categories of Spaces, Lemma
\ref{spaces-perfect-lemma-pseudo-coherent}).
Hence by Descent on Spaces, Proposition
\ref{spaces-descent-proposition-fpqc-descent-quasi-coherent}
we see that
$$
H^0(T, \mathcal{F}) = \Ker(
\prod H^0(T_i, h_i^*\mathcal{F}) \to
\prod H^0(T_i \times_T T_j, (T_i \times_T T_j \to T)^*\mathcal{F}))
$$
In this way we see that the rules in (1) and (2) satisfy
the sheaf property for fppf coverings. This means we may apply
Bootstrap, Lemma \ref{bootstrap-lemma-locally-algebraic-space-finite-type}
to see it suffices to prove the representability \'etale locally on $B$.
Moreover, we may check whether the end result is affine and
of finite presentation \'etale locally on $B$, see
Morphisms of Spaces, Lemmas \ref{spaces-morphisms-lemma-affine-local} and
\ref{spaces-morphisms-lemma-finite-presentation-local}.
Hence we may assume that $B$ is an affine scheme.

\medskip\noindent
Assume $B = \Spec(A)$ is an affine scheme. By the results of
Derived Categories of Spaces, Lemmas
\ref{spaces-perfect-lemma-pseudo-coherent},
\ref{spaces-perfect-lemma-derived-quasi-coherent-small-etale-site}, and
\ref{spaces-perfect-lemma-descend-pseudo-coherent}
we deduce that in the rest of the proof we may think of $K$ as a perfect
object of the derived category of complexes of modules on $B$
in the Zariski topology. By 
Derived Categories of Schemes, Lemmas
\ref{perfect-lemma-pseudo-coherent},
\ref{perfect-lemma-affine-compare-bounded}, and
\ref{perfect-lemma-pseudo-coherent-affine} we can find a pseudo-coherent
complex $M^\bullet$ of $A$-modules such that $K$ is the corresponding
object of $D(\mathcal{O}_B)$. Our assumption on pullbacks implies
that $M^\bullet \otimes^\mathbf{L}_A \kappa(\mathfrak p)$
has vanishing $H^{-1}$ for all primes $\mathfrak p \subset A$.
By More on Algebra, Lemma \ref{more-algebra-lemma-cut-complex-in-two}
we can write
$$
M^\bullet =
\tau_{\geq 0}M^\bullet \oplus \tau_{\leq - 1}M^\bullet
$$
with $\tau_{\geq 0}M^\bullet$ perfect with Tor amplitude in $[0, b]$
for some $b \geq 0$ (here we also have used
More on Algebra, Lemmas \ref{more-algebra-lemma-glue-perfect} and
\ref{more-algebra-lemma-glue-tor-amplitude}).
Note that in case (2) we also see that $\tau_{\leq - 1}M^\bullet = 0$
in $D(A)$ whence $M^\bullet$ and $K$ are perfect with
tor amplitude in $[0, b]$. For any $B$-scheme $g : T \to B$ we have
$$
H^0(T, H^0(Lg^*K)) = H^0(T, H^0(Lg^*\tau_{\geq 0}K))
$$
(by the dual of Derived Categories, Lemma
\ref{derived-lemma-negative-vanishing})
hence we may replace $K$ by $\tau_{\geq 0}K$ and correspondingly
$M^\bullet$ by $\tau_{\geq 0}M^\bullet$. In other words, we may
assume $M^\bullet$ has tor amplitude in $[0, b]$.

\medskip\noindent
Assume $M^\bullet$ has tor amplitude in $[0, b]$.
We may assume $M^\bullet$ is a bounded above complex of finite free
$A$-modules (by our definition of pseudo-coherent complexes, see
More on Algebra, Definition \ref{more-algebra-definition-pseudo-coherent}
and the discussion following the definition).
By More on Algebra, Lemma \ref{more-algebra-lemma-last-one-flat}
we see that $M = \Coker(M^{- 1} \to M^0)$ is flat. By
Algebra, Lemma \ref{algebra-lemma-finite-projective} we see that $M$
is finite locally free. Hence $M^\bullet$ is quasi-isomorphic to
$$
M \to M^1 \to M^2 \to \ldots \to M^d \to 0 \ldots
$$
Note that this is a K-flat complex
(Cohomology, Lemma \ref{cohomology-lemma-bounded-flat-K-flat}),
hence derived pullback of $K$ via a morphism $T \to B$ is computed
by the complex
$$
g^*\widetilde{M} \to g^*\widetilde{M^1} \to \ldots
$$
Thus it suffices to show that the functor
$$
(g : T \to B) \longmapsto
\Ker(
\Gamma(T,g^*\widetilde{M})
\to
\Gamma(T, g^*(\widetilde{M^1})
)
$$
is representable by an affine scheme of finite presentation over $B$.

\medskip\noindent
We may still replace $B$ by the members of an affine open covering
in order to prove this last statement. Hence we may assume that $M$
is finite free (recall that $M^1$ is finite free to begin with).
Write $M = A^{\oplus n}$ and $M^1 = A^{\oplus m}$. Let the map
$M \to M^1$ be given by the $m \times n$ matrix $(a_{ij})$ with
coefficients in $A$. Then $\widetilde{M} = \mathcal{O}_B^{\oplus n}$
and $\widetilde{M^1} = \mathcal{O}_B^{\oplus m}$. Thus the functor
above is equal to the functor
$$
(g : T \to B) \longmapsto
\{(f_1, \ldots, f_n) \in \Gamma(T, \mathcal{O}_T) \mid
\sum g^\sharp(a_{ij})f_i = 0,\ j = 1, \ldots, m\}
$$
Clearly this is representable by the affine scheme
$$
\Spec\left(A[x_1, \ldots, x_n]/(\sum a_{ij}x_i; j = 1, \ldots, m)\right)
$$
and the lemma has been proved.
\end{proof}
